\section{Conclusion}
\label{sec:conclusion}

We crossed knowledge with incentive on matched trivia questions for \llama and scored lie detectors and hallucination detectors on the resulting false answers.
False answers are almost entirely epistemic unless the model is told to lie: an incentive alone adds about one knowing misreport per hundred known questions, while an explicit instruction makes lies 40\% of false answers.
The standard deception probe is dominated by the prompt at any fixed threshold.
Within a prompt it responds more to knowing lies than to hallucinations (0.80), so it is not a general falsehood detector.
Hallucination detectors are closer to general detectors of falsehood and low confidence: they fire on both kinds of false answer.
Lies and hallucinations are linearly separable (0.92), mostly through a signal of question knowledge that is present before the answer.

Several steps follow from these findings.
First, we plan to repeat the design on a larger model that lies under incentive alone, so that the lie cell is not instructed.
Second, calibrating the probe threshold per prompt, or subtracting the prompt-token score, would test whether the within-prompt signal survives as a usable detector.
Third, a condition that asks for a different but still correct answer, such as an alias or paraphrase, would separate ``departs from the default answer'' from ``contradicts the belief''.
Finally, a follow-up-token probe, NLI-based semantic entropy, and human validation of a sample of lie labels would tighten the detector set and the labels.
